% Source: d34_population_balance_mechanism.md, Sections 8--10, precision-mlps 7a12986.
\section{Joint feature acquisition: conditions and full proof}
\label{app:feature-theory}

The result controls an evolving tanh network in the ordinary Euclidean GD
metric. A scalar comparison bounds the collective size of slopes, hidden
biases, and readouts; a global remainder inequality converts that bound
into output error. Its coefficients can be supplied as structural upper
bounds or evaluated along a trajectory. The latter use is a conditional
explanation, not a forecast from the initial parameters alone.

\subsection{Network, measure, and coarse compensation}
Let $\nu$ be a probability measure on $[-1,1]$ with positive input variance,
and write $\|\cdot\|=\|\cdot\|_{L^2(\nu)}$ throughout this appendix.
For the training measure, this is the sample RMS norm $\|\cdot\|_{2,Z}$.
Consider
\begin{align}
 q_\theta(x)&=d+\sum_{j=1}^W c_j\tanh(\gamma_jx+b_j),&
 L(\theta)&=\tfrac12\|q_\theta-f\|^2,\notag\\
 \theta_{n+1}&=\theta_n-\eta_n\nabla L(\theta_n),& \eta_n&\ge0.
 \label{eq:feature-network}
\end{align}
The target lies in $L^2(\nu)$ with $\|f\|>0$, and $[a]_+=\max(a,0)$.
Let $P_C$ project onto affine functions,
$P_H=I-P_C$, $f_\perp=P_Hf$, $q_H=P_Hq_\theta$, and
$e_C=P_C(q_\theta-f)$, $e_H=q_H-f_\perp$.
Jacobians $J_C=D_\theta(P_Cq_\theta)$ and $J_H=D_\theta q_H$ map the
Euclidean parameter space into these output spaces. Their adjoints use
the stated output measure; empirical adjoints therefore include the
sample normalization.

At the iterates under consideration assume $K_C=J_CJ_C^*$ is invertible
on the two-dimensional affine space. Define
\begin{align}
 \Pi&=I-J_C^*K_C^{-1}J_C,&
 \ell&=K_C^{-1}J_CJ_H^*e_H,\notag\\
 F&=\Pi J_H^*e_H,& R&=J_C^*(e_C+\ell).
 \label{eq:feature-split}
\end{align}
Then $\nabla L=F+R$, $\Pi$ is an orthogonal projection, and $J_CF=0$.
Thus $F$ includes the compensation needed to preserve affine output to
first order. Small tracking $R$ does not mean that compensation vanishes,
or that $e_C$ is exactly zero. The proof uses $\|\Pi\|\le1$;
it needs no quantitative lower singular-value bound on $J_C$.
At a singular coarse Gram matrix this formulation is not invoked.

Write $\vartheta=(\gamma,b,c)$ for the hidden parameter blocks, omitting
the output bias, and set
\begin{equation}
 r_j=(\gamma_j^2+b_j^2+c_j^2)^{1/2},\qquad
 S_k=\sum_j r_j^k,\qquad s=\|\vartheta\|_2=\sqrt{S_2}>0.
 \label{eq:feature-radius}
\end{equation}
The slope and center parameters are free to evolve. Here $b_j$ is the
hidden intercept; a center, when $\gamma_j\ne0$, is $-b_j/\gamma_j$.
No bound on center displacement is inferred by dividing by a small slope.

\subsection{Global tanh and derivative remainder bounds}
Put $u_j=\gamma_jx+b_j$ and define
\begin{equation}
 \Psi_3=-\tfrac13\sum_j c_ju_j^3,\qquad
 \Psi_5=\tfrac2{15}\sum_j c_ju_j^5,\qquad
 J_k=D_\theta(P_H\Psi_k),\quad G_k=J_k^*f_\perp\quad(k=3,5).
 \label{eq:feature-polynomials}
\end{equation}
Set $t_3=1/3$, $t_5=2/15$, $t_7=17/315$ and
\begin{equation}
 A_k=t_k\,\frac{2^{k/2}k^{k/2}}{(k+1)^{(k+1)/2}},\qquad
 C_k=(k+1)A_k,\qquad k=3,5,7.
 \label{eq:feature-constants}
\end{equation}
In particular $A_3=\sqrt6/8$. For every real $u$, the remainders after
the linear, cubic, and quintic tanh polynomials satisfy
\begin{equation}
 |\mathcal R_k(u)|\le t_k|u|^k,\qquad
 |\mathcal R_k'(u)|\le kt_k|u|^{k-1},\qquad k=3,5,7.
 \label{eq:feature-scalar-remainders}
\end{equation}
Here $\mathcal R_3=\tanh u-u$, $\mathcal R_5=\tanh u-u+u^3/3$,
and $\mathcal R_7=\tanh u-u+u^3/3-2u^5/15$.

\paragraph{Proof of the scalar bounds.}
Since $|\tanh u|\le|u|$, $|\mathcal R_3'|=\tanh^2u\le u^2$;
integration gives the cubic bound. Next,
$\mathcal R_5'=u^2-\tanh^2u$, whose magnitude is at most
$|u-\tanh u|\,|u+\tanh u|\le2|u|^4/3$; integrate again.
For the last pair, $\sup_{u\in\mathbb R}|\tanh^{(7)}u|\le272$.
To check this constant explicitly, write $z=\tanh^2u\in[0,1]$:
\[
 \tanh^{(7)}u=-272+3968z-12096z^2+13440z^3-5040z^4.
\]
The degree-four Bernstein coefficients on $[0,1/2]$ and $[1/2,1]$ are
$(-272,224,216,124,53)$ and $(53,-18,-68,8,0)$, respectively.
The Bernstein basis is nonnegative and sums to one, proving the bound.
Taylor's integral remainder for tanh through degree six and its derivative
through degree five gives $272/7!=17/315$ and $272/6!=17/45$.
The intervening coefficients vanish by parity.

\paragraph{From neurons to output and Jacobians.}
On $[-1,1]$, $|u_j|\le\sqrt2(\gamma_j^2+b_j^2)^{1/2}$.
At fixed $r_j$, maximizing
$t_k2^{k/2}|c_j|(\gamma_j^2+b_j^2)^{k/2}$ gives $A_kr_j^{k+1}$.
The squared norm of that neuron's three derivative columns is bounded
pointwise by
\[
 t_k^22^k\left[k^2c_j^2(\gamma_j^2+b_j^2)^{k-1}
                    +(\gamma_j^2+b_j^2)^k\right]\le C_k^2r_j^{2k}.
\]
The maximum occurs at $(\gamma_j^2+b_j^2)/r_j^2=k/(k+1)$.
Sum output bounds and squared column bounds, integrate against the
probability measure, and use contractivity of $P_H$. With the
Hilbert--Schmidt norm bounding the operator norm, this proves
\begin{align}
 \|q_H-P_H(\Psi_3+\Psi_5)\|&\le A_7S_8,\notag\\
 \|J_H-J_3-J_5\|_{\rm HS}&\le C_7\sqrt{S_{14}},\notag\\
 \|J_H^*f_\perp-G_3-G_5\|_2&\le C_7\|f_\perp\|\sqrt{S_{14}}.
 \label{eq:feature-remainders}
\end{align}
The same argument gives $\|q_H\|\le A_3S_4$ and
$\|J_H\|_{\rm HS}\le C_3\sqrt{S_6}$.
These are global inequalities for the exact tanh network: no neuronwise
small-argument hypothesis or polynomial replacement of training is used.

\subsection{Structural coefficients and target dependence}
For $s>0$ define
\begin{align}
 \chi_k&=W^{k/2-1}S_k/s^k,\notag\\
 \alpha_k&=W^{(k-1)/2}\|P_H\Psi_k\|/s^{k+1},\notag\\
 \beta_k&=W^{(k-1)/2}\|J_k\|_{\rm HS}/s^k,\qquad
 \zeta_k=W^{(k-1)/2}\|G_k\|_2/s^k,\quad k=3,5.
 \label{eq:feature-shapes}
\end{align}
These coefficients describe concentration, polynomial output, sensitivity,
and target coupling. They are invariant under common positive rescaling
of $\vartheta$; bounding them does not itself bound $s$.
For equal neuron norms, $\chi_k=1$; for the moments $k>2$ used here,
concentration in fewer neurons increases it.
Let $\tau_k=\|P_{\mathcal P_k}f_\perp\|$, where $\mathcal P_k$ is the
space of polynomials of degree at most $k$ for the same measure. Then
\begin{equation}
 \zeta_k\le\beta_k\tau_k,\qquad
 \alpha_k\le A_k\chi_{k+1},\qquad
 \beta_k\le C_k\sqrt{\chi_{2k}},\qquad k=3,5.
 \label{eq:feature-target-loading}
\end{equation}
Each column of $J_k$ is a polynomial of degree at most $k$, so
$J_k^*f_\perp=J_k^*P_{\mathcal P_k}f_\perp$. The other inequalities
follow from the fixed-radius bounds just proved. Thus cubic target
orthogonality makes $G_3=0$, and quintic orthogonality also makes $G_5=0$.
The mixed-sine and localized targets generally have nonzero low-degree
coupling. The numerical evaluations retain $\alpha_k,\beta_k,\zeta_k$;
substituting their concentration-only bounds gives a different, potentially
looser evaluation, whose tightness is not claimed here.

At each GD iterate $n$, take finite nonnegative upper bounds
$\widehat\alpha_{k,n},\widehat\beta_{k,n},\widehat\zeta_{k,n}$,
$\widehat\chi_{8,n},\widehat\chi_{14,n}$ on these coefficients, and
$\widehat r_n\ge\|R_n\|_2$. Exact coefficients are allowed.
For a candidate radius $B\ge0$, define
\begin{align}
 \mathcal Q_n(B)&=
 \widehat\alpha_{3,n}\frac{B^4}{W}
 +\widehat\alpha_{5,n}\frac{B^6}{W^2}
 +A_7\widehat\chi_{8,n}\frac{B^8}{W^3},\notag\\
 \mathcal J_n(B)&=
 \widehat\beta_{3,n}\frac{B^3}{W}
 +\widehat\beta_{5,n}\frac{B^5}{W^2}
 +C_7\sqrt{\widehat\chi_{14,n}}\frac{B^7}{W^3},\notag\\
 \mathcal T_n(B)&=
 \widehat\zeta_{3,n}\frac{B^3}{W}
 +\widehat\zeta_{5,n}\frac{B^5}{W^2}
 +C_7\|f_\perp\|\sqrt{\widehat\chi_{14,n}}\frac{B^7}{W^3},\notag\\
 \mathcal V_n(B)&=\mathcal J_n(B)\mathcal Q_n(B)+\mathcal T_n(B).
 \label{eq:feature-envelopes}
\end{align}
For $s_n\le B$, these bound $\|q_{H,n}\|$, $\|J_{H,n}\|$,
$\|J_{H,n}^*f_\perp\|_2$, and $\|F_n\|_2$, respectively. For the last,
\begin{equation}
 \|F_n\|_2
 \le\|J_{H,n}\|\,\|q_{H,n}\|+\|J_{H,n}^*f_\perp\|_2
 \le\mathcal V_n(B).
 \label{eq:feature-force}
\end{equation}
All four polynomials are nondecreasing in $B$. Their coefficients may evolve;
neither the Jacobian nor the target coupling is frozen.

\subsection{Precise GD theorem and proof}
\begin{theorem}[Population growth and output error for joint GD]
\label{thm:feature-precise}
Consider \eqref{eq:feature-network} from an arbitrary restart, indexed by
$n=0$. Suppose $K_{C,n}$ is invertible, $s_n>0$, and the structural
upper bounds above hold for every iterate $0\le n\le N$.
Define
\begin{equation}
 B_0=s_0,\qquad
 B_{n+1}=B_n+\eta_n[\mathcal V_n(B_n)+\widehat r_n].
 \label{eq:feature-recurrence}
\end{equation}
Whenever these quantities are finite, for $0\le n\le N$,
\begin{align}
 s_n&\le B_n,\qquad
 \sum_{k<n}\|\vartheta_{k+1}-\vartheta_k\|_2\le B_n-B_0,
 \label{eq:feature-travel}\\
 \frac{\|f-q_n\|}{\|f\|}
 &\ge\frac{[\|f_\perp\|-\mathcal Q_n(B_n)]_+}{\|f\|},\qquad
 \overline\lambda_n\le\lambda_{{\rm RMS},n}\le\frac{hB_n}{\sqrt W}.
 \label{eq:feature-consequences}
\end{align}
Here $h>0$ is the reference spacing,
$\lambda_{{\rm RMS},n}=h\|\gamma_n\|_2/\sqrt W$, and
$\overline\lambda_n=hW^{-1}\sum_j|\gamma_{j,n}|$.
\end{theorem}

\begin{proof}
Initially $s_0=B_0$. If $s_n\le B_n$, the exact GD update and
\eqref{eq:feature-force} give
\[
 \|\vartheta_{n+1}-\vartheta_n\|_2
 \le\eta_n\|F_n+R_n\|_2
 \le\eta_n[\mathcal V_n(B_n)+\widehat r_n]=B_{n+1}-B_n.
\]
The triangle inequality implies $s_{n+1}\le B_{n+1}$; induction and summation
prove \eqref{eq:feature-travel}. Orthogonal projection and the reverse
triangle inequality yield
\[
 \|q_n-f\|\ge\|q_{H,n}-f_\perp\|
 \ge[\|f_\perp\|-\|q_{H,n}\|]_+
 \ge[\|f_\perp\|-\mathcal Q_n(B_n)]_+.
\]
Finally Cauchy--Schwarz and $\|\gamma_n\|_2\le s_n$ prove both bandwidth
bounds. These arguments use the discrete updates directly, with no
continuous-time approximation or unmeasured finite-step correction.
\end{proof}

The result holds for any nonnegative step sizes for which its premises
hold. Large steps can make the comparison uninformative. Small post-transient
tracking makes it useful but is not silently set to zero. A large value
of $B_n$ or a zero output floor signals loss of this sufficient explanation;
it does not establish successful acquisition. Coefficient bounds checked only
at stored times do not certify their values at every intervening GD step.
For a pooled mean across equal-sized seeds, the per-seed bounds average;
the theorem does not interpret neuron spread as uncertainty across seeds.

\subsection{Mechanistic interpretation and conditional width scaling}
For gradient flow the same proof uses
$\dot B=\mathcal V_t(B)+\widehat r(t)$, $B(0)=s(0)$.
Locally bounded measurable coefficients make the right side locally Lipschitz
and nondecreasing in $B$. Since $D^+s\le\mathcal V_t(s)+\widehat r(t)$,
scalar comparison proves the radius, travel, and output bounds until the
comparison ceases to be finite. Effective flow sets $R=0$; full flow retains it.

Suppose $s(0)=O(1)$ and the structural coefficients and target norm have
width-independent bounds while $s$ stays in a fixed bounded interval.
Then $\mathcal Q=O(W^{-1})$, $\mathcal J=O(W^{-1})$, and the leading
target term in $\mathcal T$ is $O(W^{-1})$.
With tracking accumulation below a fixed fraction of the radius change,
an order-one increase of radius requires $\Omega(W)$ flow time.
If $f_\perp\perp\mathcal P_3$, the cubic target term vanishes,
$\mathcal V=O(W^{-2})$, and the analogous time is $\Omega(W^2)$.
The same reasoning applies to cumulative GD time $\sum_n\eta_n$.
It does not establish width-independent conditions empirically or a
universal update-count scaling under width-dependent learning rates.
Higher target gaps alone do not improve this particular $W^2$ conclusion:
the unsigned product $\mathcal J\mathcal Q$ can remain $O(W^{-2})$.

The signed balance identity provides complementary interpretation.
For $\Delta=\frac12\sum_j(\gamma_j^2+b_j^2-c_j^2)$, set
\[
 k_\theta=\sum_jc_j(\tanh u_j-u_j\operatorname{sech}^2u_j),\qquad
 E_\theta=\sum_jc_j(3\tanh u_j-u_j\operatorname{sech}^2u_j-2u_j).
\]
Since $Dq_\theta[\nabla\Delta]=-k_\theta$ and
$P_Hk_\theta=-2q_H+P_HE_\theta$, effective flow satisfies exactly
\begin{equation}
 \dot\Delta=-2\|q_H\|^2+2\langle f_\perp,q_H\rangle
 +\langle e_H,P_HE_\theta\rangle-\langle\ell,P_Ck_\theta\rangle.
 \label{eq:feature-balance}
\end{equation}
Full flow adds $\langle e_C+\ell,P_Ck_\theta\rangle$.
The negative term corrects network-generated non-affine output; target
coupling can drive positive growth. Neither the net sign nor slope decay
is asserted. Theorem~\ref{thm:feature-precise} uses the norm comparison
\eqref{eq:feature-force}, not an assumption that this signed correction
dominates. Controlling readouts together with geometry is essential:
small slopes alone do not imply the output limitation proved here.

\section{Evidence and interpretation for Section 3.5}
\label{app:feature-evidence}

\subsection{Five-million-update observation and feature assays}
The primary target is
$[\sin(2\pi x)+\frac12\sin(6\pi x)+\frac14\sin(14\pi x)]/\sqrt{21/32}$
on $[-1,1]$. Total width is 512, including 22 halo nodes per side;
the reference spacing is $h=2/467$. Full-batch FP64 training uses 2,048
midpoints and five paired initializations. One recipe per joint optimizer
is selected by median final error on 4,096 validation midpoints, requiring
finite endpoints for every seed. The 8,192-point evaluation grid is a
previously inspected resolution check, not an untouched test set.
The selected joint Adam and GD recipes use initial rates 0.05 and 0.2,
respectively, with cosine decay over all five million updates.
The supplied-feature Adam comparison uses $\lambda=1/4$, zero readout
initialization, and the selected full-horizon cosine rate $10^{-4}$.
Both constant and cosine candidates are retained in the study records.

Median joint training errors are $1.8143227\times10^{-3}$ for Adam and
$0.24393137$ for GD. Supplied-feature Adam reaches $6.6245807\times10^{-7}$.
The corresponding dense-grid values are $1.8148250\times10^{-3}$,
$0.24393487$, and $6.6198225\times10^{-7}$.
Joint Adam's median endpoint error falls about 28\% between the two- and
five-million-update experiments. This is a precision gap within the stated
budget, not permanent failure.
The pooled endpoint mean bandwidths are 0.0168845 for Adam and 0.00135795
for GD. Their neuron population standard deviations are 0.0541204 and
0.00703476. The supplied bandwidth is a comparison geometry, not a
target-independent necessary threshold. The revised three-panel Figure 4
is external to this review copy; no new width-scaling result is asserted here.

Freezing the five-million-update Adam features and training new readouts
for two million additional updates gives median dense-grid error
$1.8147369\times10^{-3}$. SVD readout fits at relative cutoffs
$10^{-12}$ and $10^{-14}$ also yield only small improvements; lowering the
cutoff to $10^{-16}$ exposes sensitivity to floating-point cancellation.
These assays support a limitation of the learned dictionary at the checked
numerical precision, not an exact-arithmetic approximation floor.
Multiplying both slope and intercept by 4 or 16 preserves each learned
center. After the same readout-training assay, median dense-grid errors
are 0.0916060 and 0.138935, respectively. Thus simple slope amplification
does not repair this geometry. These interventions change nonnested feature
spaces; they do not establish that larger slopes always help.
The assay's additional optimization budget
is separate from the end-to-end five-million-update comparison.

\subsection{Conditional theorem evaluation on evolving GD populations}
The mechanism study uses a separate cohort: total width 705, reference
spacing $h=1/256$, 2,048 uniform midpoint samples, and restarts after
20,000 ordinary-GD updates at step size 0.002. Each continuation runs
another 100,000 updates, or 200 units of flow time.
Seed 30 covers six targets: an orthogonal degree-five mixture, mixed sine,
a shifted Gaussian, a compact bump, a smooth step, and an absolute-value kink.
The mixture is $0.3P_0+0.4P_1+\sqrt{0.75}P_5$, with $P_k$ orthonormal
for the empirical measure. The other unnormalized targets are
\[
 \begin{gathered}
 \sin(2\pi x)+\tfrac12\sin(6\pi x)+\tfrac14\sin(14\pi x),\\
 e^{-((x+0.35)/0.22)^2},\qquad \tanh(14(x-0.31)),\qquad |x+0.23|,\\
 f_{\rm bump}(x)=
 \begin{cases}
 \exp\!\left(1-\dfrac1{1-u^2}\right),& |u|<1,\quad u=(x-0.35)/0.22,\\
 0,&|u|\ge1.
 \end{cases}
 \end{gathered}
\]
Each is divided by its RMS on the original training grid.
Four additional seed-31 restarts cover step, kink, Gaussian, and bump;
they were selected for weak initial comparison margins before their
continuation outcomes were inspected.

The discrete scalar recurrence uses structural diagnostics sampled every
unit of flow time, interpolating its nonnegative polynomial coefficients
and tracking norm to native GD steps. This implements a numerical evaluation
of the conditional theorem; interpolation is not a proved upper enclosure
between stored samples. In the six principal GD runs the minimum sampled
relative-error floors are, in the target order above,
$0.865734$, $0.912194$, $0.784451$, $0.879212$, $0.479360$, and $0.426952$.
The corresponding executed endpoint errors are
$0.866025$, $0.912564$, $0.788073$, $0.881027$, $0.489675$, and $0.427513$.
The minimum floors and endpoint errors are distinct summaries, not an
every-step tightness statistic. All ten growing comparisons remain useful
over the full continuation, with seed-31 minimum floors from 0.426695 to
0.878467. Principal endpoint RMS-bandwidth bounds range from
$3.63\times10^{-4}$ to $5.18\times10^{-4}$, versus observed values from
$2.09\times10^{-4}$ to $2.51\times10^{-4}$.

The saved verification includes independent automatic derivatives, global
remainder checks, exact discrete moment identities, and a solvable
positive-growth comparison. Step refinement changes population size and
normalized slope RMS by less than $4.8\times10^{-7}$ relatively across both
seeds. Decimating the structural diagnostics changes the final GD comparison
radius by less than $9\times10^{-6}$ relatively. These checks support
numerical resolution; they do not supply interval-arithmetic certificates
for the conditional premises.

The theorem explains this post-transient regime through evolving population
structure. The older width-177 archive includes regimes where large moments
make the polynomial comparison uninformative. No coverage claim for all
23 archived targets, all training ages, or the full five-million-update
Figure 4 trajectories follows. Adam's coordinate-dependent scaling and
momentum also fall outside the ordinary-GD proof. Its observed precision
gap motivates the same empirical question without establishing the same
quantitative mechanism.
